\begin{problem}[Recursive sequence and limiting value]
Let $(x_n)$ be defined by
\[
x_1=16,\qquad x_{n+1}=2+\frac{1}{x_n}\quad (n\ge 1).
\]
\begin{enumerate}[(i)]
\item Find exact values of $x_2,x_3,x_4$.
\item Assuming $x_n\to L>0$, show that $L=1+\sqrt{2}$.
\item Prove that
\[
|x_{n+1}-L|=\frac{|x_n-L|}{Lx_n},
\]
and deduce that for $n\ge 2$,
\[
|x_{n+1}-L|<\frac14|x_n-L|.
\]
\item Hence show $x_n\to 1+\sqrt{2}$.
\end{enumerate}
\end{problem}

\begin{solution}
\textbf{(i)}
\[
x_2=2+\frac1{16}=\frac{33}{16},\quad
x_3=2+\frac{16}{33}=\frac{82}{33},\quad
x_4=2+\frac{33}{82}=\frac{197}{82}.
\]
For $n\ge 2$, we have $x_n>2$ because each term is $2+\frac{1}{x_{n-1}}$ with positive second part.

\textbf{(ii)} If $x_n\to L$, then taking limits in the recurrence gives
\[
L=2+\frac1L \implies L^2-2L-1=0.
\]
So $L=1\pm\sqrt2$. Since terms are positive and $>2$ from $n\ge2$, we take
\[
L=1+\sqrt2.
\]

\textbf{(iii)}
\[
x_{n+1}-L=\left(2+\frac1{x_n}\right)-\left(2+\frac1L\right)
=\frac{L-x_n}{Lx_n},
\]
thus
\[
|x_{n+1}-L|=\frac{|x_n-L|}{Lx_n}.
\]
Now $L=1+\sqrt2>2$ and $x_n>2$ for $n\ge2$, so $Lx_n>4$. Hence
\[
|x_{n+1}-L|<\frac14|x_n-L|,\quad n\ge2.
\]

\textbf{(iv)} Repeatedly applying the inequality gives geometric decay of the error:
\[
|x_{n}-L|\le \left(\frac14\right)^{n-2}|x_2-L|\to 0.
\]
Therefore $x_n\to L=1+\sqrt2$.
\end{solution}

\begin{takeaways}
\item For recursive limits, first solve the fixed-point equation.
\item A contraction inequality gives a rigorous convergence proof.
\item The same method works for many rational recursions.
\end{takeaways}

\begin{problem}[The Dynamics of Iteration]
Consider the function
\[
g(x)=\sqrt{x+2}.
\]
We define an iterative sequence by
\[
x_{n+1}=g(x_n),\qquad x_1=1.
\]

\begin{enumerate}[(i)]
\item Calculate $x_2,x_3,$ and $x_4$ to three decimal places. By solving
\[
L=\sqrt{L+2},
\]
determine the exact positive limit $L$ that the sequence appears to be approaching.
\item Show that
\[
0<g'(x)<\frac12
\]
for all $x>0$.
\item Prove that
\[
|x_{n+1}-L|<\frac12|x_n-L|.
\]
Hence deduce that the sequence converges to $L$ for any initial value $x_1>0$.
\item Define a function $f(x)$ such that Newton's method would converge to the same limit $L$ found in part (i), and write down the specific iterative formula for this $f(x)$.
\end{enumerate}
\end{problem}
\newpage
\begin{hint}
\begin{itemize}
\item \textbf{For (i):} Substitute repeatedly into $x_{n+1}=\sqrt{x_n+2}$, then solve the fixed-point equation.
\item \textbf{For (ii):} Differentiate $g(x)=\sqrt{x+2}$ and note where the derivative is largest for $x>0$.
\item \textbf{For (iii):} Use either the mean value theorem or rationalise $\sqrt{x_n+2}-\sqrt{L+2}$.
\item \textbf{For (iv):} Rearrange $L=\sqrt{L+2}$ into a polynomial equation.
\end{itemize}
\end{hint}
\begin{solution}
\textbf{(i)}
\[
x_2=\sqrt3\approx1.732,\quad
x_3=\sqrt{2+\sqrt3}\approx1.932,\quad
x_4\approx1.983.
\]
If $x_n\to L>0$, then
\[
L=\sqrt{L+2}\implies L^2-L-2=0
\implies (L-2)(L+1)=0.
\]
Thus the positive limit is
\[
L=2.
\]

\textbf{(ii)}
\[
g'(x)=\frac{1}{2\sqrt{x+2}}.
\]
For $x>0$, this is positive and
\[
g'(x)<\frac{1}{2\sqrt2}<\frac12.
\]

\textbf{(iii)} Since $L=2$ and $g(L)=L$,
\[
\begin{aligned}
|x_{n+1}-L|
&=\left|\sqrt{x_n+2}-\sqrt{L+2}\right|\\
&=\frac{|x_n-L|}{\sqrt{x_n+2}+\sqrt{L+2}}\\
&<\frac12|x_n-L|,
\end{aligned}
\]
because $x_n>0$ and $\sqrt{L+2}=2$. Repeating this inequality gives
\[
|x_n-L|<\left(\frac12\right)^{n-1}|x_1-L|\to0,
\]
so $x_n\to L=2$ for any $x_1>0$.

\textbf{(iv)} From the fixed-point equation,
\[
L^2-L-2=0.
\]
So take
\[
f(x)=x^2-x-2.
\]
Newton's method gives
\[
x_{n+1}=x_n-\frac{f(x_n)}{f'(x_n)}
=x_n-\frac{x_n^2-x_n-2}{2x_n-1}.
\]
\end{solution}
\begin{takeaways}
\item Fixed points often reveal the likely limit of an iterative sequence.
\item A contraction inequality proves convergence by forcing errors to shrink geometrically.
\item Newton's method can target the same limit by solving the fixed-point equation as a root problem.
\end{takeaways}

\begin{problem}[The Geometry of Harmonic Means]
A sequence $(h_n)$ is in harmonic progression (HP) if the sequence of reciprocals
\[
x_n=\frac{1}{h_n}
\]
forms an arithmetic progression (AP).

\begin{enumerate}[(i)]
\item Given that the third term of an HP is $1$ and the sixth term is $\frac12$, find the first term $h_1$ and the common difference $d$ of the underlying AP.
\item The harmonic mean $H$ of two positive numbers $a$ and $b$ is defined so that $a,H,b$ is an HP. Show algebraically that
\[
H=\frac{2ab}{a+b}.
\]
\item Let
\[
A=\frac{a+b}{2},\qquad G=\sqrt{ab}
\]
be the arithmetic and geometric means of $a$ and $b$. Prove that $G^2=AH$. Hence prove that, for distinct positive real numbers $a$ and $b$,
\[
H<G<A.
\]
\item Consider the infinite harmonic series
\[
\sum_{n=1}^{\infty}\frac1n.
\]
By comparing the sum of the first $2^k$ terms to blocks of constant lower bounds, prove that this harmonic series does not have a finite limiting sum.
\end{enumerate}
\end{problem}
\newpage
\begin{hint}
\begin{itemize}
\item \textbf{For (i):} Convert the HP terms into AP terms $x_3$ and $x_6$, then use $x_n=a+(n-1)d$.
\item \textbf{For (ii):} If $a,H,b$ is an HP, then $\frac1a,\frac1H,\frac1b$ is an AP.
\item \textbf{For (iii):} Multiply $A$ and $H$. For the inequality, use $(a-b)^2>0$ or $(\sqrt a-\sqrt b)^2>0$.
\item \textbf{For (iv):} Group terms as $\left(\frac13+\frac14\right)$, then $\left(\frac15+\cdots+\frac18\right)$, and so on.
\end{itemize}
\end{hint}
\begin{solution}
\textbf{(i)} Since $x_n=\frac1{h_n}$,
\[
x_3=1,\qquad x_6=2.
\]
For the underlying AP,
\[
x_6-x_3=3d \implies 2-1=3d \implies d=\frac13.
\]
Also
\[
x_1=x_3-2d=1-\frac23=\frac13,
\]
so
\[
h_1=\frac1{x_1}=3.
\]

\textbf{(ii)} Since $\frac1a,\frac1H,\frac1b$ is an AP, the middle term is the average of the outer terms:
\[
\frac1H=\frac12\left(\frac1a+\frac1b\right)=\frac{a+b}{2ab}.
\]
Therefore
\[
H=\frac{2ab}{a+b}.
\]

\textbf{(iii)} Using the formulas for $A$ and $H$,
\[
AH=\frac{a+b}{2}\cdot\frac{2ab}{a+b}=ab=G^2.
\]
For $a\ne b$,
\[
(\sqrt a-\sqrt b)^2>0
\implies a+b>2\sqrt{ab}
\implies A>G.
\]
Since $G^2=AH$ and all quantities are positive, $A>G$ gives
\[
H=\frac{G^2}{A}<G.
\]
Hence
\[
H<G<A.
\]

\textbf{(iv)} Group the harmonic series into dyadic blocks:
\[
1+\frac12+\left(\frac13+\frac14\right)+\left(\frac15+\cdots+\frac18\right)+\cdots.
\]
Each block after the first has sum greater than or equal to $\frac12$:
\[
\frac13+\frac14>\frac14+\frac14=\frac12,\qquad
\frac15+\cdots+\frac18>4\cdot\frac18=\frac12,
\]
and similarly
\[
\frac1{2^{j-1}+1}+\cdots+\frac1{2^j}>2^{j-1}\cdot\frac1{2^j}=\frac12.
\]
Thus the partial sums exceed
\[
1+\frac12+\underbrace{\frac12+\frac12+\cdots+\frac12}_{\text{arbitrarily many terms}},
\]
so they grow without bound. Therefore the harmonic series has no finite limiting sum.
\end{solution}
\begin{takeaways}
\item Harmonic progressions are arithmetic progressions viewed through reciprocals.
\item For positive distinct numbers, the classical means satisfy $H<G<A$.
\item A sequence can have terms tending to zero while its infinite series still diverges.
\end{takeaways}

\begin{problem}[Chebyshev composition]

The first-kind Chebyshev polynomials $T_n(x)$ are defined by the recurrence
\[
T_0(x)=1,\quad T_1(x)=x,\quad T_{n+1}(x)=2xT_n(x)-T_{n-1}(x).
\]

It is given that $T_n(\cos\theta)=\cos(n\theta)$ for all $n\ge0$ and $\theta\in\R$. (Do NOT prove this fact.)

For first-kind Chebyshev polynomials, prove that
\[
T_m(T_n(x))=T_{mn}(x)
\]
for all $x$ of the form $x=\cos\theta$. Then use this result to find
\[
T_2(T_3(x)).
\]
\end{problem}
\begin{hint}
Use $T_n(\cos\theta)=\cos(n\theta)$. For the explicit polynomial, use $T_6(x)$.
\end{hint}
\begin{solution}
Let $x=\cos\theta$. Then
\[
T_m(T_n(x))=T_m(\cos(n\theta))=\cos(mn\theta)=T_{mn}(\cos\theta)=T_{mn}(x).
\]
Therefore
\[
T_2(T_3(x))=T_6(x).
\]
Using the recurrence after $T_5(x)=16x^5-20x^3+5x$,
\[
T_6(x)=2xT_5(x)-T_4(x)=32x^6-48x^4+18x^2-1.
\]
\end{solution}
\begin{takeaways}
\item The trigonometric definition makes Chebyshev composition simple.
\item Composition of first-kind Chebyshev polynomials multiplies the indices.
\item Hard-looking polynomial composition can often be avoided by using structure.
\end{takeaways}
